Can a multitape Turing machine transpose a $d\times d$ matrix of $\lceil\log_2 d\rceil$-bit entries in time linear in its
length~$L$? We show that an oblivious machine that does so must route a linear amount of information through
\emph{macroscopic crossings}: pairs of revisits of one tape, each taking the head a fixed fraction of~$L$ away from its cell,
that interleave as a queue's would. Precisely, for every set of at least $\alpha L$ output bits, every set of steps that separates
their output steps from the earlier ends of such crossings in the time--space dependency graph has at least $c$ times as many
elements as the output set. Here $c$ and the fraction depend only on the number of tapes, the alphabet, the number of states, the
time constant and~$\alpha$. The expansion can be found on a single coarsely monotone leg of a single tape, in a coarse skeleton of the
run drawn from a finite family independent of~$d$, and at least two work heads must make revisits of linear excursion. Behind these
statements is a dichotomy for arbitrary linear-time traces: either such crossings are reached with linear expansion, or a separator
exists that is cheaper than the acyclic set it leaves in the cone relation, which no transposer can have. For adaptive machines the
conclusions hold on every trace followed by an exponentially small but fixed fraction of the inputs, and on the heavy traces of a bit
simulator, with the same tapes and alphabet, of every linear-time transposer; carrying them back to the machine's own traces is open.
We also prove that every cut between a set of input rows and the outputs they feed into a set of output rows carries the corresponding
information whatever the coding, and we exclude classes of machines with boundedly many macroscopic phases, oblivious and adaptive.
Explicit linear-time schedules, realized by machines that compute other functions, meet these conditions in every test we could run
(numerically, at $d\le128$), so the conditions are not known to decide the question. We show that one set of steps separates the
inputs from the outputs for every pair of row sets at once, so counting such cuts cannot by itself go beyond linear time, and we state
the multi-scale question that remains. The main theorems are proved in Lean~4; one proposition and one corollary are proved in the
text, and the exclusions of individual schedule families are given as proof sketches. The constants of the main theorems are explicit
and astronomically small; those of the two-heads theorem are polynomial, and those of one class exclusion come from a compactness
argument.
